% sections/09-validity.tex — §9: validity as first-class measurement (spec §6 protocol).
\section{Evaluation Validity}
\label{sec:validity}

Our third insight treats validity as a measurable property of evaluations, not a limitations
paragraph. We organize the evidence as a progression --- information leakage, then dataset-level
invalidity, then oracle manipulation --- because each tier changes what a reported number can
mean.

\subsection{V2/V3: Information Leakage, Measured by Intervention}
The cleanest datum in our corpus is controlled. Fang et al.\ measure GPT-4-based agents on 15
real one-day vulnerabilities under two information conditions: with the CVE description
supplied, the agent autonomously exploited 87\%; without it, success fell to 7\%, while every
baseline (GPT-3.5, eight open-source models, ZAP, Metasploit) scored zero in both
conditions~\cite{fang2024agents}. The twelve-fold swing is produced by a single
supplied-knowledge variable --- the CVE text functions as a partial answer key --- so any
capability claim inherits its information condition. We therefore report such results only with
their conditions attached, and the same discipline applies to temporal leakage: every
vulnerability in that benchmark was publicly disclosed before model training cutoffs.

\subsection{V1/V5: Dataset-Level Invalidity, Measured by Reconstruction}
PrimeVul demonstrates that the detection literature's benchmark layer was structurally
unsound: label noise, near-duplicate pairs across splits, and pre-split public exposure
inflated reported performance to the point where paired, chronologically split evaluation
collapses strong-model accuracy toward chance~\cite{primevul2024}. Risse and B\"ohme reach a
compatible verdict from measurement theory, arguing that popular vulnerability datasets score
models on the wrong exam~\cite{rissebohme2024}. The corpus also contains a live saturation
illustration: Argusee reports 100\% on CyberSecEval~2 buffer-overflow cases while its own
documentation requires human-supplied entry points~\cite{darknavyargusee2025} --- ceiling
scores coexisting with bounded autonomy. Agentic CTF suites inherit a further V1 exposure:
their tasks derive from publicly written-up competitions, and no included work quantifies the
inflation this induces; we return to it as roadmap experiment E2.

\subsection{V4/V5: Oracle Manipulation and Environment Brittleness, Documented Operationally}
AIxCC supplied operational evidence on outcome definitions. PoV re-execution establishes
behavior only for the tested proof; teams documented mitigation-gaming patches --- recompiling
with MTE/PAC or wrapping entry points in broad \texttt{catch} blocks --- that pass PoVs without
semantic repair, prompting optional functional tests as secondary oracles and an accuracy
modifier that made patch correctness the binding score term~\cite{taafc2025}. Environment
leakage is exemplified by the winning CRS nearly disabling itself when organizers prefixed
challenge directories with ``ossfuzz,'' breaking a ``fuzz''-substring heuristic --- load-bearing
benchmark-specific logic invisible until the environment shifted~\cite{taafc2025}. On the
production side, leaderboard-style counting conflates heterogeneous outcomes: of XBOW's
${\sim}1{,}060$ submissions, roughly 39\% resolved into duplicate or informative categories
rather than new/resolved findings --- categories that reflect discovery overlap and program
policy as much as agent error, making rank non-decomposable without category-aware
accounting~\cite{xbowtop12025}.

\subsection{Administrative Record Correction: A Separate Category}
DARPA's closeout corrects its own synthetic-vulnerability count from 70 to 63 with an editor's
note~\cite{darparesults2025}, and prize accounting requires three scoped figures rather than
one (\$29.5M cumulative commitment, \$8.5M final pool, \$30.5M
distributed~\cite{darpascoring2025,openssfaixcc2026}). These are administrative record
corrections, not measurement-validity failures, but they carry a methodological moral we adopt
as policy: even authoritative primary sources require version- and date-aware verification,
and press intermediaries inherit whatever instability the primary contains.

\subsection{A Contamination-Control Protocol}
As a concrete deliverable, we specify the minimal protocol any agentic capability evaluation
should publish: (i)~\emph{temporal firewall} --- freeze the task set at date $T$, evaluate only
models whose training cutoffs precede $T-\Delta$ for disclosed $\Delta$, and release the task
set only after evaluation; (ii)~\emph{provenance screening} --- machine-check every task against
public write-up corpora and NVD text up to $T$, reporting overlap rates rather than asserting
cleanliness; (iii)~\emph{information-condition disclosure} --- state exactly which artifacts
(descriptions, patches, stack traces, credentials) are supplied to the agent, Fang-style;
(iv)~\emph{category-aware outcomes} --- report validated/exploitable/duplicate/informative as
separate counts with adjudication rules fixed before running; and (v)~\emph{adjudication
audit} --- a second scorer re-labels a random 10\% sample with agreement statistics. No
included evaluation satisfies all five today; E3 (Sec.~\ref{sec:roadmap}) operationalizes the
gap.

\subsection{What Survives}
Validity constraints do not render the corpus uninterpretable. Condition-controlled results
survive (the 7\% no-description condition), third-party adjudication survives (Ada Logics'
reproduction of 27 AIxCC-derived issues~\cite{openssfaixcc2026}), and platform triage provides
weak external checks on vendor claims. What does not survive is unqualified headline numbers ---
which are precisely the ones that travel furthest.
